% ===========================================================================
% Shared notation for the internal note and the compendium.
%
% Conventions follow Mezard, "Generative diffusion updated notes" (Dec 2025),
% sections 2-3.1.1, and the statistical-physics usage of the surrounding
% literature. Where that source and the project's earlier code disagreed, the
% source wins and the code is translated in the text rather than the reverse.
%
% Two deliberate departures from the source, both noted where they occur:
%
%   1. The source uses n for the number of data samples and d for the ambient
%      dimension. Here the ambient dimension IS the chain length and is the
%      quantity every graphical-model statement is indexed by, so it is written
%      n, and the number of observed sequences is M. A translation sentence
%      appears at first use.
%
%   2. The source writes the decay factor as e^{-t} throughout. We keep e^{-t}
%      in displayed equations and never abbreviate it to alpha_t.
% ===========================================================================

% --- Diffusion -------------------------------------------------------------
\newcommand{\Pz}{P_0}                       % clean distribution
\newcommand{\Pt}{P_t}                       % noised distribution at time t
\newcommand{\Dt}{\Delta_t}                  % 1 - e^{-2t}
\newcommand{\score}{S}                      % S(x,t) = grad log P_t(x)
% Standardised Gaussian white noise: <eta(t)eta(t')> = delta(t-t'). The source's section 2
% uses this convention in its Eq (3) but a variance-2 convention in its Eqs (1) and (18); we
% fix the standardised one everywhere and carry the sqrt(2) explicitly in both the forward
% and the reverse SDE, so that Delta_t = 1 - e^{-2t} is consistent throughout.
\newcommand{\gwn}{\eta}

% Posterior mean, in the source's bracket notation rather than E[a|x].
\newcommand{\pmean}[1]{\langle #1 \rangle_{x,t}}
\newcommand{\pmeansite}[1]{\langle a_{#1} \rangle_{x,t}}

% --- Chain and graphical model --------------------------------------------
\newcommand{\nsites}{n}                     % chain length = ambient dimension
\newcommand{\nseq}{M}                       % number of observed sequences
\newcommand{\ngrid}{K}                      % quadrature grid points
\newcommand{\kernel}{W}                     % transition kernel W(a'|a)
\newcommand{\like}{\ell_t}                  % unary likelihood factor
\newcommand{\lmsg}{\overrightarrow{m}}      % forward (left) message
\newcommand{\rmsg}{\overleftarrow{m}}       % backward (right) message
\newcommand{\innov}{\varphi}                % innovation density
\newcommand{\corr}{\rho}                    % autoregression coefficient
\newcommand{\ivar}{q}                       % innovation variance

% --- Estimation ------------------------------------------------------------
\newcommand{\params}{\theta}
\newcommand{\Qfun}{\mathcal{Q}}             % EM auxiliary function
\newcommand{\Xis}{\Xi}                      % expected transition statistic
\newcommand{\Lik}{\mathcal{L}}              % marginal log-likelihood

% --- Operators -------------------------------------------------------------
\newcommand{\E}{\mathbb{E}}
\newcommand{\Prob}{\mathbb{P}}
\newcommand{\Reals}{\mathbb{R}}
\newcommand{\Cov}{\operatorname{Cov}}
\newcommand{\Var}{\operatorname{Var}}
\newcommand{\KL}{D_{\mathrm{KL}}}
\newcommand{\normal}[2]{\mathcal{N}\!\left(#1,#2\right)}

% --- Estimators ------------------------------------------------------------
\newcommand{\mmse}{\widehat{a}_{\mathrm{MMSE}}}
\newcommand{\lmmse}{\widehat{a}_{\mathrm{LMMSE}}}

% --- Repository pointers (shared by the note and the compendium) ------------
\providecommand{\coderef}[2]{\texttt{\small #1}\,$\rightarrow$\,\texttt{\small #2}}
\providecommand{\codefile}[1]{\texttt{\small #1}}
